% ============================================================
% คำนำ (Preface)
% ============================================================
\chapter*{คำนำ}
\addcontentsline{toc}{chapter}{คำนำ}

ระบบสื่อสารไร้สายในแปลงเกษตรกรรมสมัยใหม่ถือเป็นโครงสร้างพื้นฐานที่มีบทบาทสำคัญอย่างยิ่งต่อความแม่นยำในการเก็บรวบรวมข้อมูลโทรมาตรและการสั่งการอุปกรณ์ควบคุมอัตโนมัติ สภาพแวดล้อมทางการเกษตรในภูมิภาคเขตร้อนชื้น เช่น สวนผลไม้เศรษฐกิจและโรงเรือนปลูกพืช มักเผชิญกับอุปสรรคทางกายภาพจากการลดทอนของคลื่นวิทยุโดยทรงพุ่มใบพืชที่มีปริมาณน้ำสูง ตลอดจนระยะทางระหว่างโหนดตรวจวัดที่กระจายตัวในพื้นที่กว้าง เทคโนโลยีโครงข่ายแบบเมชตามมาตรฐาน Zigbee 3.0 บนพื้นฐานข้อกำหนด IEEE 802.15.4 จึงเป็นคำตอบทางวิศวกรรมที่เหมาะสมที่สุด ด้วยคุณสมบัติเด่นด้านการจัดเส้นทางเครือข่ายด้วยตนเอง การซ่อมแซมเส้นทางส่งสัญญาณแบบพลวัต และการประหยัดพลังงานอย่างสูงสุด

คู่มือวิชาการฉบับนี้เรียบเรียงขึ้นเพื่อเป็นเอกสารหลักในการศึกษา วิจัย และลงมือปฏิบัติการพัฒนาระบบสื่อสารไร้สาย Zigbee 3.0 สำหรับงานวิศวกรรมฟิสิกส์เกษตรแม่นยำ โดยมุ่งเน้นการถ่ายทอดองค์ความรู้ตั้งแต่ระดับรากฐานทางฟิสิกส์คลื่นวิทยุ การวิเคราะห์เปรียบเทียบเชิงสมรรถนะร่วมกับเทคโนโลยีอื่น การประยุกต์ใช้งานโหนดเซนเซอร์ตรวจวัดสภาพดิน สภาพบรรยากาศ และรังสีดวงอาทิตย์ การเชื่อมต่อกับหัววัดมาตรฐานอุตสาหกรรมผ่านสะพานสื่อสาร RS485 Modbus RTU ไปจนถึงการติดตั้งและบริหารจัดการเกตเวย์โครงข่ายร่วมกับ Zigbee2MQTT เพื่อบูรณาการข้อมูลสู่แพลตฟอร์มคลาวด์และระบบควบคุมฟาร์มอัจฉริยะ ตลอดจนการพัฒนาเฟิร์มแวร์บนไมโครคอนโทรลเลอร์รุ่นใหม่ ESP32-C6 สถาปัตยกรรม RISC-V

ผู้เรียบเรียงหวังเป็นอย่างยิ่งว่า คู่มือเล่มนี้จะเป็นประโยชน์ต่อนักศึกษา นักวิจัย วิศวกร ตลอดจนผู้สนใจในเทคโนโลยีอินเทอร์เน็ตของสรรพสิ่งทางการเกษตร เพื่อนำไปประยุกต์ใช้ในการยกระดับผลิตภาพทางการเกษตรและขับเคลื่อนนวัตกรรมเทคโนโลยีเกษตรกรรมดิจิทัลของประเทศให้ก้าวหน้าอย่างมั่นคงและยั่งยืน

\vspace{1.5cm}
\begin{flushright}
    \textbf{ผู้ช่วยศาสตราจารย์ ดร.ชีวะ ทัศนา}\\
    กลุ่มวิจัยฟิสิกส์เกษตรดิจิทัลและปัญญาประดิษฐ์ฝังตัว\\
    คณะวิทยาศาสตร์และเทคโนโลยี มหาวิทยาลัยราชภัฏรำไพพรรณี\\
    กันยายน 2570
\end{flushright}
\clearpage
